\scanpage{559}
\mspage{470}
On a l'impression que la raison \unsure{pour} laquelle
on a de tels contre-exemples, c'est que
l'on a des $\unsure{\leftrightarrows|\rightleftarrows}$-images essentielles qui ne sont
\deleted{\unsure{satisfont à}} $L_\pi$-stables dans $M$ pour aucun
cardinal $\pi\in\unsure{\mathcal U|U}$. Cela amène à \unsure{modifier}
la conjecture ainsi : soit $\pi$ un cardinal
tel que $M$ \unsure{satisfasse|satisfait} à $L_\pi$, et soit
$\underline S(\pi)$ la plus petite sous-catégorie \unsure{str. pl|str. ple} de
$M$, $L_\pi$-stable dans $M$, et qui contienne
l'image ess. de $f$. (On pourrait appeler
$\underline S(\pi)$ la $L_\pi$-image essentielle.) Alors $\underline S(\pi)$
est accessible dans $M$. De plus, la famille
(décroissante) des $\underline S(\pi)$, pour $\pi$ variable,
est stationnaire, i.e il existe un cardinal
$\pi$ tel que $\pi'\geq\pi$ implique $\underline S(\pi')=\underline S(\pi)$.
Il est clair \inserted{alors} que $\underline S(\pi)$ est \unsure{la|le} plus \unsure{petite|petit}
sous-catégorie de $M$ qui soit accessible dans $M$,
et qui contienne l'image essentielle de
$f$. La même question se pose pour
un foncteur $g^{\circ}:N^{\circ}\longrightarrow M$, où $g:N\mathrel{\unsure{-|\longrightarrow}}M^{\circ}$
est un foncteur accessible, $N$ et $M$ accessibles.
Je conjecture que la réponse est encore
\unsure{affirmative}.
% archive sheet number: 558
